% ============================================================
%   Hardware-Aware SE Networks for ECG Anomaly Detection on FPGAs
%   arXiv preprint — IEEE-compatible single-column format
% ============================================================

\documentclass[12pt, letterpaper]{article}

% --- Layout ---
\usepackage[top=1in, bottom=1in, left=1.25in, right=1.25in]{geometry}
\usepackage{setspace}
\setstretch{1.15}

% --- Fonts and typography ---
\usepackage{times}
\usepackage{microtype}

% --- Mathematics ---
\usepackage{amsmath}
\usepackage{amssymb}

% --- Graphics and figures ---
\usepackage{graphicx}
\usepackage{caption}
\usepackage{subcaption}
\usepackage{float}
\captionsetup{font=small, labelfont=bf, skip=6pt}

% --- Tables ---
\usepackage{booktabs}
\usepackage{tabularx}
\usepackage{multirow}
\usepackage{array}

% --- Lists ---
\usepackage{enumitem}

% --- Units ---
\usepackage{siunitx}
\sisetup{detect-all}

% --- Colors ---
\usepackage{xcolor}
\usepackage[breaklinks=true]{hyperref}

% Citations(IEEE-style numbered)
\usepackage{cite}

% --- Hyperlinks ---

\usepackage{url}

% --- Section formatting ---
\usepackage{titlesec}
\titleformat{\section}
  {\normalfont\large\bfseries}{\Roman{section}.}{0.5em}{}
\titleformat{\subsection}
  {\normalfont\normalsize\bfseries}{\Alph{subsection}.}{0.5em}{}
\titlespacing*{\section}{0pt}{1.2ex plus 0.5ex}{0.8ex}
\titlespacing*{\subsection}{0pt}{1.0ex plus 0.3ex}{0.5ex}

% --- Abstract formatting ---
\usepackage{abstract}
\renewcommand{\abstractname}{\textbf{Abstract}}
\setlength{\absleftindent}{0pt}
\setlength{\absrightindent}{0pt}

% --- Keywords command ---
\newcommand{\keywords}[1]{%
  \vspace{0.4em}%
  
\noindent\textbf{Index Terms---}#1%
  \vspace{0.6em}%
}

% --- Paragraph spacing ---
\setlength{\parskip}{3pt}
\setlength{\parindent}{1.2em}

% ============================================================

\begin{document}

% --- Title block ---
\title{%
  \textbf{Hardware-Efficient Squeeze-and-Excitation Networks for Real-Time ECG Anomaly Detection Using hls4ml on FPGAs}%
}

\author{%
  Nitesh Kumar\\
  \normalsize Department of Electrical Engineering\\
  \normalsize National Institute of Technology Patna\\
  \normalsize Patna, Bihar 800005, India\\
  \small\texttt{niteshk.ug23.ee@nitp.ac.in}
  \and
  Ashwani Kumar\\
  \normalsize Department of Electrical Engineering\\
  \normalsize National Institute of Technology Patna\\
  \normalsize Patna, Bihar 800005, India\\
  \small\texttt{ashwani@nitp.ac.in}
}

\date{} % suppress date for arXiv preprint style

\maketitle

% ============================================================
\begin{abstract}
% ============================================================
We present a hardware-efficient Squeeze and Excitation (SE)
channel-attention network for cardiac anomaly detection on
field programmable gate arrays (FPGAs), designed 
for direct deployment via the open-source \texttt{hls4ml} high-level synthesis framework. The proposed 1,674-parameter architecture applies SE channel
attention on the post-pooled temporal representation, a design choice
that eliminates tensor-broadcasting operations and ensures correct HLS
code generation from Keras model graphs.
Trained on the ECG5000
benchmark (500 training samples, 4,500 test samples, 140 time steps per
sample), the best-performing software model achieves an AUC-ROC of
93.32\% and an accuracy of 91.18\%.
Through a systematic five-point fixed-point
precision sweep from 8-bit to 24-bit representations, we identify
\texttt{ap\_fixed\allowbreak<24,12>} as the minimum viable precision for
this architecture;
coarser configurations degrade the AUC to near-chance
levels. After fixed-point conversion at \texttt{ap\_fixed<24,12>}, the
hardware C-simulation model records an AUC-ROC of 92.84\% with a logic
fidelity of 99.9\% relative to the floating-point reference (one
disagreement per 1,000 stratified test samples).
Vivado RTL synthesis on
the Xilinx xc7z020clg400-1 yields a per-sample inference latency of
2.0\,\si{\micro\second} at 100\,MHz, corresponding to a throughput
approximately 142,000 times faster than the real-time requirement for
continuous 500\,Hz ECG monitoring.
We also document and provide a
reproducible fix for a softmax normalization defect in
\texttt{hls4ml}~1.3.0 C-simulation that silently corrupts multi-class
probability outputs.
\end{abstract}

\keywords{Anomaly detection; channel attention; \texttt{hls4ml}; FPGA; fixed-point quantization; Squeeze and Excitation networks;  cardiac monitoring; edge inference.}

\vspace{1em}
\hrule


% ============================================================
\section{Introduction}
\label{sec:intro}
% ============================================================
Cardiovascular diseases account for approximately 17.3 million deaths
annually~\cite{Roth2020}, and the ECG remains the primary non-invasive
tool for diagnosing cardiac arrhythmias.
The proliferation of wearable
biosensors has shifted ECG acquisition to continuous ambulatory settings,
generating data volumes that exceed the capacity of expert
review~\cite{Hannun2019} and creating an acute demand for automated,
real-time classification at the point of recording.
Deep convolutional networks have achieved cardiologist-level arrhythmia
detection on large annotated corpora~\cite{Rajpurkar2017, Hannun2019},
yet these models are designed for data-center hardware and are
incompatible with the milliwatt power budgets, constrained memory, and
sub-millisecond latency requirements of embedded cardiac monitors.
FPGAs
address this through custom parallel datapaths that realize full
neural-network inference at power well below that of a
GPU~\cite{Guo2019}, while retaining field-updateability for model
revision without hardware re-fabrication.
The \texttt{hls4ml}
framework~\cite{Duarte2018, Fahim2021} lowers the implementation barrier
by automatically converting Keras layer graphs to Vivado HLS C++ with
configurable fixed-point types (\texttt{ap\_fixed}) and reuse factors
that control loop-unrolling depth.
Attention mechanisms improve temporal classification accuracy, but
Transformer-style self-attention~\cite{Vaswani2017} scales as
$\mathcal{O}(N^2)$ with sequence length: a 140-sample ECG window requires
$140 \times 140 = 19{,}600$ inner products per attention head, straining
the DSP resources and on-chip memory of embedded FPGAs.
The
Squeeze-and-Excitation (SE) channel-attention mechanism~\cite{Hu2018}
achieves $\mathcal{O}(1)$ complexity with respect to sequence length by
replacing token-to-token scoring with global average pooling and
channel-wise excitation on fixed-size descriptors.
However, naïve SE
integration within a temporal encoder introduces tensor-broadcast operations and attention weights of shape $(1,C)$ multiplied against
feature maps of shape $(L,C)$ , that \texttt{hls4ml}~1.3.0 does not
model correctly, producing silent inference errors in C-simulation.
This paper presents a hardware-aware SE-attention network for ECG anomaly
detection, designed for reliable \texttt{hls4ml} conversion.
Positioning
SE attention after global average pooling reduces all attention operations
to element-wise multiplications on 16-dimensional vectors, eliminating
broadcasting from the generated HLS code.
A softmax normalization defect
in \texttt{hls4ml}~1.3.0 is additionally identified and resolved. The
best software model achieves AUC-ROC 93.32\% on ECG5000;
at
\texttt{ap\_fixed<24,12>}, the hardware model achieves AUC-ROC 92.84\%
with 99.9\% logic fidelity and 2.0\,\si{\micro\second} inference latency
at 100\,MHz on the Xilinx xc7z020clg400-1. \\

\medskip
\noindent\textbf{Contributions.}\;
This paper makes the following
specific contributions.

\begin{itemize}[leftmargin=2em, noitemsep, topsep=4pt]

  \item \textbf{Hardware-compatible SE attention.}
    Relocating the SE block to post-pooled 16-dimensional vectors
    eliminates all tensor-broadcast operations from the
    \texttt{hls4ml}-generated HLS code, enabling correct RTL synthesis
    without toolchain modification.
\item \textbf{Softmax normalization fix.}
    We document a reproducible C-simulation defect in
    \texttt{hls4ml}~1.3.0 in which both binary class probabilities
    simultaneously exceed~0.5, and provide a toolchain-independent remedy
    applicable to any multi-class \texttt{hls4ml} workflow.
\item \textbf{Fixed-point precision characterization.}
    A systematic five-configuration sweep from \texttt{ap\_fixed<8,4>}
    to \texttt{ap\_fixed<24,12>} establishes \texttt{ap\_fixed<24,12>}
    as the minimum viable precision for 140-step ECG inference;
coarser
    formats produce degenerate or near-random predictions.
\item \textbf{End-to-end FPGA validation.}
    Hardware AUC-ROC 92.84\%, logic fidelity 99.9\%, and
    Vivado-synthesized inference latency 2.0\,\si{\micro\second} at
    100\,MHz on the Xilinx xc7z020clg400-1, approximately
    142,000$\times$ faster than the real-time requirement for 500\,Hz
    ECG monitoring.
\end{itemize}

% ============================================================
\section{Related Work}
\label{sec:related}
% ============================================================

\subsection{ECG Classification with Deep Learning}

Rajpurkar et al.~\cite{Rajpurkar2017} demonstrated that a 34-layer
residual CNN trained on 64,000 single-lead recordings matches cardiologist
sensitivity across 12 arrhythmia classes;
Hannun et al.~\cite{Hannun2019}
extended this to ambulatory consumer-wearable data, achieving ten-class
accuracy exceeding a panel of six cardiologists.
Across the UCR benchmark
suite, residual convolutional architectures consistently outperform
classical time-series classifiers~\cite{IsmailFawaz2019}, with a shallow
residual baseline~\cite{Wang2017} remaining competitive against more
complex models.
Kachuee et al.~\cite{Kachuee2018} showed that
representations pre-trained on MIT-BIH Arrhythmia Database recordings
transfer effectively to ECG5000~\cite{Dau2019} without dataset-specific
feature engineering.
These high-accuracy models assume server-class
hardware and are ill-suited to the resource constraints of embedded
cardiac monitors.
\subsection{Attention Mechanisms and Squeeze-and-Excitation Networks}

Transformer self-attention~\cite{Vaswani2017} captures long-range
temporal dependencies but scales as $\mathcal{O}(N^2)$ in both memory
and computation, which is prohibitive for FPGAs with limited DSP slices
and on-chip BRAM.
The SE mechanism~\cite{Hu2018} offers a
hardware-efficient alternative: global average pooling produces a
fixed-size channel descriptor independent of sequence length, and a
compact bottleneck network with sigmoid gating generates multiplicative
channel re-weighting at $\mathcal{O}(1)$ cost relative to temporal
extent.
Batch normalization~\cite{Ioffe2015} and residual
connections~\cite{He2016} are incorporated as standard complements to
stabilize optimization and preserve gradient flow.
The critical
implementation constraint, addressed in this work, is that SE excitation
must operate on post-pooled fixed-dimensional vectors to avoid
shape-variable broadcast operations that are incompatible with static HLS
synthesis.
\subsection{FPGA-Based Neural Network Inference and hls4ml}

FPGA neural-network accelerator architectures are surveyed
in~\cite{Guo2019}.
FINN~\cite{Umuroglu2017} and
FINN-R~\cite{Blott2018} demonstrated high-throughput streaming dataflow
designs exploiting bitwise and ternary computation, at the cost of
restricting the representable model class.
The \texttt{hls4ml}
framework~\cite{Duarte2018, Fahim2021} preserves network topology,
translating Keras or PyTorch layer graphs to Xilinx HLS C++ via
automated layer-by-layer conversion with configurable fixed-point
arithmetic and loop-unrolling depth.
Aarrestad et
al.~\cite{Aarrestad2021} showed that \texttt{hls4ml} converted 1D
convolutional networks achieve latency and resource utilization
competitive with hand-crafted HLS on UltraScale devices;
Coelho et
al.~\cite{Coelho2021} extended the framework to heterogeneous per-layer
mixed-precision assignment. These studies address particle-physics trigger
systems operating at nanosecond timescales.
To the best of our knowledge, no prior work has combined SE channel attention with
\texttt{hls4ml} synthesis and a systematic multi-precision fixed-point
characterization for ECG classification on an embedded Zynq-class FPGA.

\section{Methodology}
\label{sec:method}

\subsection{Dataset}

The ECG5000 benchmark~\cite{Dau2019} contains 5{,}000 single-channel
heartbeat segments of $L = 140$ uniformly sampled time steps each.
The standard split assigns 500 samples to training and 4{,}500 to
testing, with a class ratio of 58.4\% Normal (class~0) to 41.6\%
Anomaly (class~1) maintained in both partitions.
Figures~\ref{fig:ecg_samples} and~\ref{fig:class_dist} show
representative waveforms and the class balance, respectively.

\begin{figure}[!ht]
  \centering
  \includegraphics[width=1\linewidth]{fig_ecg_samples.png}
  \caption{Normal (top) and anomalous (bottom) ECG segments from ECG5000}
  
  \label{fig:ecg_samples}
\end{figure}

\begin{figure}[!ht]
  \centering
  \includegraphics[width=0.65\linewidth]{fig_class_distribution.png}
  \caption{ECG5000 class distribution across training and test partitions.}
  \label{fig:class_dist}
\end{figure}

% --------------------------------
\subsection{Proposed Architecture}
\label{sec:arch}
% --------------------------------

Figure~\ref{fig:architecture} illustrates the SE-ECG model, and
Table~\ref{tab:arch} provides the layer-by-layer summary.
The network processes the $(140, 1)$ input through a temporal encoder,
a global pooling step, and a channel-wise SE attention block before
producing two-class logit scores.
\begin{figure}[!ht]
  \centering
  \includegraphics[width=0.75\linewidth]{Full_Architecture.png}
  \caption{SE-ECG architecture (1{,}674 parameters).
The softmax layer
    is excluded from the FPGA sub-model.}
  \label{fig:architecture}
\end{figure}

\begin{table}[!ht]
  \centering
  \caption{Architecture Summary}
  \label{tab:arch}
  \small
  \renewcommand{\arraystretch}{1.1}
  \begin{tabular}{llr}
    \toprule
    \textbf{Layer} & \textbf{Output Shape} & \textbf{Params} \\
    \midrule
    Input                          & $(140,\;1)$   &     0 \\
    Dense(16), time-dist.\ + BN + ReLU & $(140,\;16)$ &    96 \\
    Dense(32), time-dist.\ + BN + ReLU & $(140,\;32)$ 
&   672 \\
    Dense(16), time-dist.\ + Add + BN  & $(140,\;16)$ &   592 \\
    GlobalAvgPool1D                & $(16,)$       &     0 \\
    Dense(8), ReLU                 & $(8,)$        &   136 \\
    Dense(16), Sigmoid     
        & $(16,)$       &   144 \\
    Multiply + Add (residual)      & $(16,)$       &     0 \\
    Dense(2) — logit output        & $(2,)$        &    34 \\
    Softmax$^\dagger$              & $(2,)$     
   &     0 \\
    \midrule
    \textbf{Total}                 &               & \textbf{1{,}674} \\
    \bottomrule
    \multicolumn{3}{l}{%
      \footnotesize$^\dagger$Excluded from FPGA sub-model.}
  \end{tabular}
\end{table}

\noindent\textbf{Temporal Encoder.}\quad
A time-distributed \texttt{Dense(16)} embedding layer applies a shared
linear projection with batch normalization~\cite{Ioffe2015} and ReLU to
each of the $L$ time steps.
A two-layer feed-forward network (FFN)
expands the channel dimension to 32 and contracts it back to 16, with
a residual skip connection~\cite{He2016} over the entire expansion-contraction
block.
A final batch normalization layer is applied to the residual sum,
yielding $\mathbf{R} \in \mathbb{R}^{L \times C}$ ($C = 16$). \\
\noindent\textbf{Global Average Pooling.}\quad
The temporal dimension is collapsed to a $C$-dimensional channel
descriptor:
\begin{equation}
  z_c = \frac{1}{L}\sum_{t=1}^{L} R_{t,c}, \quad c = 1,\ldots,C
  \label{eq:gap}
\end{equation}

\noindent\textbf{Squeeze-and-Excitation Block.}\quad
The SE mechanism~\cite{Hu2018} computes channel attention weights from
$\mathbf{z}$ through a bottleneck dense network with sigmoid gating:
\begin{equation}
  \mathbf{s} = \sigma\!\left(
    \mathbf{W}_2\;\mathrm{ReLU}\!\left(
      \mathbf{W}_1\,\mathbf{z} + \mathbf{b}_1
    \right) + \mathbf{b}_2
  \right)
  \label{eq:se}
\end{equation}
where $\mathbf{W}_1 \in \mathbb{R}^{r \times C}$,
$\mathbf{W}_2 \in \mathbb{R}^{C \times r}$, and $r = 8$.
The attended output is formed by element-wise channel reweighting
combined with a residual:
\begin{equation}
  \mathbf{h} = \mathbf{z} + \bigl(\mathbf{s} \odot \mathbf{z}\bigr)
  \label{eq:se_res}
\end{equation} \\

\noindent\textbf{Output Logits.}\quad
\begin{equation}
  \hat{\mathbf{y}} = \mathbf{W}_o\,\mathbf{h} + \mathbf{b}_o,
  \quad \mathbf{W}_o \in \mathbb{R}^{K \times C},\quad K = 2
  \label{eq:logits}
\end{equation}

A softmax layer follows \eqref{eq:logits} during training but is
deliberately implemented as a \emph{separate} Keras
\texttt{Activation} layer so that the logit sub-model can be
extracted without weight modification for FPGA conversion.

\textbf{Hardware-Compatibility Rationale.}\quad
Applying SE attention on $\mathbf{z} \in \mathbb{R}^{C}$
(post-pooling) rather than on $\mathbf{R} \in \mathbb{R}^{L \times C}$
is an intentional hardware-compatibility decision.
In the
pre-pooling variant, the attention weight vector $\mathbf{s}
\in \mathbb{R}^{C}$ must be broadcast to shape
$\mathbb{R}^{L \times C}$ before the element-wise product in
\eqref{eq:se_res}.
The \texttt{hls4ml} HLS code generator resolves
this broadcast incorrectly, computing only the first $C$ of the
required $L \times C$ products.
Post-pooling attention eliminates
this issue: every Multiply and Add in \eqref{eq:se_res} takes
operands of identical shape $(C,)$, producing correct HLS code.

% --------------------------------
\subsection{FPGA Implementation}
\label{sec:fpga}
% --------------------------------

\noindent\textbf{hls4ml workflow.}\quad
Figure~\ref{fig:hls4ml} shows the \texttt{hls4ml} workflow ~\cite{Duarte2018}.
The trained logit sub-model is converted to a Xilinx Vivado HLS C++
project via hls4ml.
The project is first compiled with \texttt{g++} for C-level simulation
and then synthesized by Vivado HLS to RTL, yielding resource utilization
and timing reports.
\begin{figure}[!ht]
  \centering
  \includegraphics[width=0.90\linewidth]{hls4ml_wrokflow.png}
  \caption{The workflow of hls4ml.~\cite{Duarte2018}}
  \label{fig:hls4ml}
\end{figure}

\noindent\textbf{Softmax Normalization Fix.}\quad
\texttt{hls4ml}~1.3.0 C-simulation produces softmax outputs in which
both class probabilities simultaneously exceed~0.5, violating
$\sum_k p_k = 1$, regardless of the softmax strategy (Standard or Stable) configured in the layer settings.
The fix is to exclude the
softmax layer from the HLS model entirely, allow the FPGA to output raw
logits $\hat{\mathbf{y}} \in \mathbb{R}^2$, and compute probabilities
via numerically stable softmax in Python:
$p_k = \exp(\hat{y}_k - \hat{y}_{\max}) /
\sum_j \exp(\hat{y}_j - \hat{y}_{\max})$.
This requires no modification to the \texttt{hls4ml} source code.\\

\noindent\textbf{Fixed-Point Precision and Configuration.}\quad
Table~\ref{tab:precision} reports hardware AUC-ROC for five
\texttt{ap\_fixed<W,I>} configurations evaluated on a 200-sample
balanced validation subset.
Configurations below 20 bits produce
degenerate or substantially degraded inference; \texttt{ap\_fixed<24,12>}
yields the highest hardware AUC (0.924) and is selected for all
subsequent evaluations.
Table~\ref{tab:fpga_cfg} lists the complete
synthesis configuration.

\begin{table}[!ht]
  \centering
  \caption{Fixed-Point Precision Sweep}
  \label{tab:precision}
  \small
  \begin{tabular}{lccl}
    \toprule
    \textbf{Precision} & \textbf{HW AUC} &
    \textbf{Sensitivity} & \textbf{Status} \\
    \midrule
    \texttt{ap\_fixed<8,4>}   & 0.500 & 0.000 & Degenerate \\
    \texttt{ap\_fixed<12,6>}  & 0.263 & 0.247 & Insufficient \\
    \texttt{ap\_fixed<16,8>}  & 0.718 & 0.659 & Functional \\
    \texttt{ap\_fixed<20,10>} & 0.922 & 0.800 & Good \\
    \textbf{\texttt{ap\_fixed<24,12>}}
        
                      & \textbf{0.924} & \textbf{0.800}
                              & \textbf{Selected} \\
    \bottomrule
  \end{tabular}
\end{table}

\begin{table}[!ht]
  \centering
  \caption{hls4ml Synthesis Configuration}
  \label{tab:fpga_cfg}
  \begin{tabular}{ll}
    \toprule
    \textbf{Parameter} & \textbf{Value} \\
    \midrule
    Target device       
  & Xilinx xc7z020clg400-1 \\
    \texttt{hls4ml} version & 1.3.0 \\
    Precision             & \texttt{ap\_fixed<24,12>} \\
    ReuseFactor           & 1 \\
    Strategy              & Latency \\
    Model converted       & Logit sub-model (no softmax) \\
    Clock target         
 & 100\,MHz \\
    \bottomrule
  \end{tabular}
\end{table} 
 \\
% --------------------------------
\subsection{Experimental Setup}
\label{sec:setup}
% --------------------------------

\noindent\textbf{Training.}\quad
The network is optimized with Adam~\cite{Kingma2015} using class-weighted
sparse categorical cross-entropy:
\begin{equation}
  \mathcal{L} = -\frac{1}{N}\sum_{i=1}^{N}\sum_{k=0}^{1}
  w_k\,y_{ik}\log p_{ik}, \quad w_k = \frac{N}{2\,N_k}
  \label{eq:loss}
\end{equation}
yielding $w_0{=}0.856$ (Normal) and $w_1{=}1.202$ (Anomaly).
To characterize sensitivity to weight initialization under the small
training set, the full pipeline is repeated with five fixed random seeds
(42, 123, 456, 789, 2024).
Table~\ref{tab:train} summarizes all
hyperparameters.  All training is performed on Google Colab with an
NVIDIA T4 GPU using TensorFlow~2.20 and \texttt{tf\_keras}.
\begin{table}[!ht]
  \centering
  \caption{Training Hyperparameters}
  \label{tab:train}
  \begin{tabular}{ll}
    \toprule
    \textbf{Hyperparameter} & \textbf{Value} \\
    \midrule
    Optimiser               & Adam \\
    Learning rate           & $5 \times 10^{-4}$ \\
    Batch size              & 32 \\
    Maximum epochs         
 & 60 \\
    LR reduction factor     & 0.5 (patience: 5 epochs) \\
    Early stopping patience & 12 epochs (val.\ loss) \\
    Class weights           & $w_0{=}0.856$,\;
$w_1{=}1.202$ \\
    Random seeds            & 42, 123, 456, 789, 2024 \\
    \bottomrule
  \end{tabular}
\end{table} \\

\noindent\textbf{Evaluation Metrics.}\quad
Software performance is reported on the full 4{,}500-sample test set for
the best seed, and as mean $\pm$ standard deviation across all five
seeds.
Hardware C-simulation is evaluated on a stratified
1{,}000-sample subset (584~Normal, 416~Anomaly) against which the
software model is also measured to enable direct comparison.
Reported metrics are AUC-ROC, accuracy, sensitivity, specificity, and F1-score.

Logic fidelity measures the prediction agreement rate between the
software and hardware models:
\begin{equation}
  \Phi = \frac{\bigl|\{i : \hat{c}_i^{\,\mathrm{SW}}
  = \hat{c}_i^{\,\mathrm{HW}}\}\bigr|}{N}
  \label{eq:fidelity}
\end{equation}
A value of $\Phi = 1$ indicates identical predictions across all samples.

% ============================================================
%  PHASE 3 — Results, Limitations, Conclusion, References
%  Append after phase2_content.tex and before
%  \bibliographystyle{IEEEtran} in main.tex
% ============================================================

% ============================================================
\section{Results and Discussion}
\label{sec:results}
% ============================================================

\subsection{Software Performance}

Figure~\ref{fig:training} shows accuracy and loss curves across the five
seeds. Four seeds converge to high validation accuracy; one seed (456)
collapses to near-zero sensitivity, highlighting the instability that can
arise with only 500 training samples and motivating multi-seed evaluation.

\begin{figure}[!ht]
  \centering
  \includegraphics[width=1.0\linewidth]{fig_training_curves.png}
  \caption{Training accuracy and loss curves across five seeds.}
  \label{fig:training}
\end{figure}

Table~\ref{tab:sw_results} reports software performance on the full
4,500-sample test set. The best-seed model (seed 42, 60 epochs)
achieves an AUC-ROC of 93.32\%, with high specificity (96.16\%)
indicating strong normal-rhythm rejection. Sensitivity (84.20\%)
reflects the inherent difficulty of detecting anomalies from 500
training samples. Across all five seeds, AUC exhibits low variance
($0.8968 \pm 0.0222$), whereas sensitivity is substantially more
variable ($0.6782 \pm 0.2191$), consistent with the
majority-class collapse observed in seed~456.

Figure~\ref{fig:sw_eval} shows the ROC curve and confusion matrix for
the best-seed model.

\begin{figure}[!ht]
  \centering
  \includegraphics[width=1.0\linewidth]{fig_sw_evaluation.png}
  \caption{ROC curve and confusion matrix for the best software model (seed 42, $N = 4{,}500$).}
  \label{fig:sw_eval}
\end{figure}

\begin{table}[!ht]
  \centering
  \caption{Software Model Performance on ECG5000 Test Set ($N = 4{,}500$)}
  \label{tab:sw_results}
  \small
  \begin{tabular}{lccccc}
    \toprule
    & \textbf{AUC} & \textbf{Acc.} & \textbf{Sens.} & \textbf{Spec.} & \textbf{F1} \\
    \midrule
    Best seed (42) & 0.9332 & 0.9118 & 0.8420 & 0.9616 & 0.8882 \\
    Mean $\pm$ Std & $0.897{\pm}0.022$ & $0.855{\pm}0.091$ & $0.678{\pm}0.219$ & $0.981{\pm}0.013$ & $0.771{\pm}0.193$ \\
    \bottomrule
  \end{tabular}
\end{table}

% ---
\subsection{Hardware Evaluation}
% ---

Table~\ref{tab:hw_compare} compares software and hardware predictions on
the 1,000-sample balanced validation subset under \texttt{ap\_fixed<24,12>}.
The hardware model matches or marginally exceeds the software model on
every metric: sensitivity is identical at 84.38\%, and the AUC difference
is $-0.03$\%. The single prediction disagreement (logic fidelity $\Phi = 0.999$,
see \eqref{eq:fidelity}) occurs on a sample whose softmax probabilities
fall within 0.012 of the 0.5 decision boundary in the software model,
indicating quantization-induced rounding at the margin rather than a
systematic class-boundary shift. Figure~\ref{fig:cm} shows confusion
matrices for both models alongside the degenerate \texttt{ap\_fixed<12,6>}
case for comparison.

Figure~\ref{fig:agreement} illustrates the prediction agreement distribution. 

\begin{table}[!ht]
  \centering
  \caption{Software vs.\ Hardware Comparison (\texttt{ap\_fixed<24,12>}, $N = 1{,}000$)}
  \label{tab:hw_compare}
  \small
  \begin{tabular}{lccc}
    \toprule
    \textbf{Metric} & \textbf{SW (float32)} & \textbf{HW (\texttt{ap\_fixed<24,12>})} & $\boldsymbol{\Delta}$ \\
    \midrule
    AUC-ROC & 0.9287 & 0.9284 & $-$0.0003 \\
    Accuracy & 0.9130 & 0.9140 & $+$0.0010 \\
    Sensitivity & 0.8438 & 0.8438 & $\phantom{+}$0.0000 \\
    Specificity & 0.9623 & 0.9640 & $+$0.0017 \\
    F1-score & 0.8897 & 0.8909 & $+$0.0012 \\
    Logic Fidelity ($\Phi$) & --- & 0.9990 & --- \\
    \bottomrule
  \end{tabular}
\end{table}

\begin{figure}[!ht]
  \centering
  \includegraphics[width=1.0\linewidth]{fig_confusion_matrices.png}
  \caption{Confusion matrices for the software model, hardware model and degenerate hardware model.}
  \label{fig:cm}
\end{figure}

\begin{figure}[!ht]
  \centering
  \includegraphics[width=0.98\linewidth]{fig_sw_hw_agreement.png}
  \caption{Software-hardware prediction agreement ($N = 1{,}000$). 999 samples produce identical class predictions.}
  \label{fig:agreement}
\end{figure}

% ---
\subsection{FPGA Synthesis Results}
% ---

Table~\ref{tab:resources} reports post-synthesis resource utilisation
for the logit sub-model on the Xilinx xc7z020clg400-1. DSP utilisation
is high (95.5\%), as expected for a 24-bit fixed-point design with
ReuseFactor~$= 1$: the fully unrolled weight accumulations map
directly to the device's DSP48E1 slices, with overflow implemented in
look-up-table-based constant multipliers. LUT utilisation (64.8\%)
reflects this overflow and the pipelined temporal datapath. BRAM usage
is minimal (1.4\%) because all 1,674 weights are embedded as
compile-time constants in the synthesised RTL rather than stored in
block RAM.

\begin{table}[!ht]
  \centering
  \caption{Vivado Synthesis Results on Xilinx xc7z020clg400-1}
  \label{tab:resources}
  \begin{tabular}{lrrr}
    \toprule
    \textbf{Resource} & \textbf{Used} & \textbf{Available} & \textbf{Utilisation} \\
    \midrule
    LUTs & 34{,}475 & 53{,}200 & 64.8\% \\
    FFs & 18{,}961 & 106{,}400 & 17.8\% \\
    DSPs & 210 & 220 & 95.5\% \\
    BRAMs & 2 & 140 & 1.4\% \\
    \midrule
    \textbf{Latency} & \multicolumn{2}{l}{197 cycles at 100\,MHz} & \textbf{2.0\,\si{\micro\second}} \\
    \bottomrule
  \end{tabular}
\end{table}

At 100\,MHz, the 197-cycle pipeline yields a per-sample inference latency
of 2.0\,\si{\micro\second}. For continuous 500\,Hz ECG monitoring, a
heartbeat window of 140 samples arrives every 280\,ms, making the
real-time throughput requirement approximately 3.6 inferences per second.
The synthesized design exceeds this by a factor of approximately
$140{,}000$, providing substantial margin for concurrent monitoring tasks
or higher-frequency sampling scenarios. For context, \texttt{hls4ml}
deployments in particle-physics trigger systems achieve latencies of
70--75\,ns on larger Kintex-7 and UltraScale+ devices~\cite{Duarte2018,Aarrestad2021};
the comparatively longer latency here is a direct consequence of the
140-step temporal pipeline inherent to cardiac sequence processing.

% ============================================================
\section{Limitations}
\label{sec:limits}
% ============================================================

Three limitations of the current work merit explicit acknowledgement.
First, the training set comprises only 500 samples, producing high
seed-to-seed variance in sensitivity (standard deviation 0.219) and
suggesting that performance estimates may not be tight; larger annotated
datasets or data augmentation could reduce this variance. Second, all
evaluations use a single benchmark (ECG5000), a 140-step subset of the
UCR archive; generalization to clinical 12-lead recordings with different
sampling rates and segment lengths remains unverified. Third, hardware
validation is confined to HLS C-simulation. Physical bitstream
deployment on the xc7z020 board and end-to-end timing measurements on
real ECG data are left as future work. Additionally, the near-saturation
of DSP resources (95.5\%) leaves minimal headroom: any expansion of the
architecture, wider channels, additional layers, or a higher
reduction ratio in the SE block would require either increasing the
ReuseFactor above~1 or targeting a device with more DSP slices.

% ============================================================
\section{Conclusion}
\label{sec:conclusion}
% ============================================================

We presented a hardware-aware Squeeze and Excitation channel-attention
network for ECG anomaly detection on FPGAs, deployed end-to-end via the
\texttt{hls4ml}  framework. The 1,674-parameter model achieves
a software AUC-ROC of 93.32\% on the ECG5000 benchmark, retaining
92.84\% AUC after 24-bit fixed-point conversion with a logic fidelity of
99.9\%. Vivado synthesis on the Xilinx xc7z020clg400-1 yields a
per-sample latency of 2.0\,\si{\micro\second}--approximately 142,000
times faster than the real-time requirement for 500\,Hz ECG monitoring.

Three findings of broader relevance to \texttt{hls4ml}-based biomedical
deployment are reported. First, positioning SE attention after global
average pooling so that all channel-wise operations act on
fixed-size one-dimensional vectors and is a necessary condition for correct
HLS code generation; pre-pooling attention introduces tensor-broadcast
operations that the current toolchain mishandles silently. Second,
\texttt{hls4ml}~1.3.0 contains a softmax normalization defect in
C-simulation that can be reproducibly resolved by converting only the
logit sub-model and applying softmax externally. Third, a systematic
precision sweep establishes \texttt{ap\_fixed<24,12>} as the minimum
viable fixed-point format for 140-step cardiac sequences; configurations
at 8 and 12 bits produce degenerate or near-random inference.

Future work will address physical FPGA deployment and timing validation
on real ECG streams, evaluation on larger multi-lead clinical datasets,
and architectural extensions such as multi-head channel attention or
hierarchical temporal pooling within the resource constraints of the
xc7z020 device.

%% ============================================================
\section*{Acknowledgements}
%% ============================================================

We thank the \texttt{hls4ml} development teams at CERN and Fermilab for maintaining an open-source framework that democratizes FPGA-based machine learning inference beyond high-energy physics applications. We also acknowledge the free GPU resources provided by Google Colab, which supported the training and synthesis experiments conducted in this work.

% ============================================================
%  BIBLIOGRAPHY
% ============================================================

\begin{thebibliography}{99}

\bibitem{Roth2020}
G.~A.~Roth \textit{et al.},
``Global burden of cardiovascular diseases and risk factors, 1990--2019: Update from the GBD 2019 study,''
\textit{J. Am. Coll. Cardiol.}, vol.~76, no.~25, pp.~2982--3021, 2020.

\bibitem{Hannun2019}
A.~Y.~Hannun \textit{et al.},
``Cardiologist-level arrhythmia detection and classification in ambulatory electrocardiograms using a deep neural network,''
\textit{Nature Medicine}, vol.~25, no.~1, pp.~65--69, 2019.

\bibitem{Rajpurkar2017}
P.~Rajpurkar, A.~Y.~Hannun, M.~Haghpanahi, C.~Bourn, and A.~Y.~Ng,
``Cardiologist-level arrhythmia detection with convolutional neural networks,''
arXiv:1707.01836, 2017.

\bibitem{Hu2018}
J.~Hu, L.~Shen, and G.~Sun,
``Squeeze-and-excitation networks,''
in \textit{Proc. IEEE Conf. Comput. Vis. Pattern Recognit. (CVPR)},
pp.~7132--7141, 2018.

\bibitem{Vaswani2017}
A.~Vaswani \textit{et al.},
``Attention is all you need,''
in \textit{Advances in Neural Information Processing Systems (NeurIPS)},
vol.~30, 2017.

\bibitem{Duarte2018}
J.~Duarte \textit{et al.},
``Fast inference of deep neural networks in FPGAs for particle physics,''
\textit{J. Instrum.}, vol.~13, no.~07, p.~P07027, 2018.

\bibitem{Fahim2021}
F.~Fahim \textit{et al.},
``hls4ml: An open-source codesign workflow for implementing machine learning in FPGAs,''
in \textit{Proc. ACM/SIGDA Int. Symp. Field-Programmable Gate Arrays (FPGA)},
2021.

\bibitem{Aarrestad2021}
T.~Aarrestad \textit{et al.},
``Fast convolutional neural networks on FPGAs with hls4ml,''
\textit{Mach. Learn.: Sci. Technol.}, vol.~2, no.~4, p.~045015, 2021.

\bibitem{Coelho2021}
C.~N.~Coelho Jr. \textit{et al.},
``Automatic heterogeneous quantization of deep neural networks for low-latency inference on the edge for particle detectors,''
\textit{Nature Machine Intelligence}, vol.~3, pp.~675--686, 2021.

\bibitem{Guo2019}
K.~Guo \textit{et al.},
``A survey of FPGA-based neural network inference accelerators,''
\textit{ACM Trans. Reconfigurable Technol. Syst.}, vol.~12, no.~1,
pp.~1--26, 2019.

\bibitem{Umuroglu2017}
Y.~Umuroglu \textit{et al.},
``FINN: A framework for fast, scalable binarized neural network inference,''
in \textit{Proc. ACM/SIGDA Int. Symp. Field-Programmable Gate Arrays (FPGA)},
pp.~65--74, 2017.

\bibitem{Blott2018}
M.~Blott \textit{et al.},
``FINN-R: An end-to-end deep-learning framework for fast exploration of quantized neural networks,''
\textit{ACM Trans. Reconfigurable Technol. Syst.},
vol.~11, no.~3, pp.~1--23, 2018.

\bibitem{IsmailFawaz2019}
H.~Ismail Fawaz, G.~Forestier, J.~Weber, L.~Idoumghar, and P.-A.~Muller,
``Deep learning for time series classification: A review,''
\textit{Data Mining Knowl. Discov.},
vol.~33, no.~4, pp.~917--963, 2019.

\bibitem{Wang2017}
Z.~Wang, W.~Yan, and T.~Oates,
``Time series classification from scratch with deep neural networks: A strong baseline,''
in \textit{Proc. Int. Joint Conf. Neural Netw. (IJCNN)},
pp.~1578--1585, 2017.

\bibitem{Kachuee2018}
M.~Kachuee, S.~Fazeli, and M.~Sarrafzadeh,
``ECG heartbeat classification: A deep transferable representation,''
in \textit{Proc. IEEE Int. Conf. Healthcare Informatics (ICHI)},
pp.~443--444, 2018.

\bibitem{Dau2019}
H.~A.~Dau \textit{et al.},
``The UCR time series archive,''
\textit{IEEE/CAA J. Automatica Sinica},
vol.~6, no.~6, pp.~1293--1305, 2019.

\bibitem{Ioffe2015}
S.~Ioffe and C.~Szegedy,
``Batch normalization: Accelerating deep network training by reducing internal covariate shift,''
in \textit{Proc. Int. Conf. Machine Learning (ICML)},
pp.~448--456, 2015.

\bibitem{He2016}
K.~He, X.~Zhang, S.~Ren, and J.~Sun,
``Deep residual learning for image recognition,''
in \textit{Proc. IEEE Conf. Comput. Vis. Pattern Recognit. (CVPR)},
pp.~770--778, 2016.

\bibitem{Kingma2015}
D.~P.~Kingma and J.~Ba,
``Adam: A method for stochastic optimization,''
in \textit{Int. Conf. Learn. Representations (ICLR)},
2015.

\end{thebibliography}

\end{document}
% ============================================================
%  END OF DOCUMENT
% ============================================================